%% 07_llm_judges.tex - LLM-as-Judge

\section{LLM-as-Judge}
\label{sec:llm-judges}

Using LLMs to evaluate other LLM outputs (``LLM-as-judge'') has gained traction as a scalable alternative to human evaluation~\citep{zheng2023judging}. This section examines when this approach works, when it fails, and provides concrete examples of effective judge prompts.

\subsection{The Appeal of LLM Judges}

LLM judges offer several advantages over human evaluation. They provide scale, enabling evaluation of thousands of outputs without recruiting annotators. They provide speed, generating evaluations in seconds rather than days. They provide cost efficiency, running evaluations at a fraction of the cost of human annotation. They provide consistency, avoiding the fatigue effects and inter-rater variability that affect human evaluators.

\begin{empirical}[MT-Bench Findings~\citep{zheng2023judging}]
In evaluating chat assistants, GPT-4 as a judge achieved \textbf{over 80\% agreement} with human preferences, matching the agreement level between human annotators themselves. This makes LLM-as-judge viable for many evaluation tasks.
\end{empirical}

LLM judges work well for relative comparisons between two outputs rather than absolute quality assessment. They excel on tasks where quality dimensions are well-defined and unambiguous. They are valuable for initial screening before human evaluation and for high-volume regression testing where human evaluation is impractical.

\subsection{Known Biases and Failure Modes}

LLM judges exhibit several documented biases that evaluators must account for. Table~\ref{tab:llm-biases} summarizes the key biases with their typical magnitude and mitigation strategies.

\begin{table}[h]
\centering
\footnotesize
\caption{Documented biases in LLM-as-judge evaluation with mitigation strategies.}
\label{tab:llm-biases}
\begin{tabular}{@{}p{2.2cm}p{4.5cm}p{2cm}p{3cm}@{}}
\toprule
\textbf{Bias} & \textbf{Description} & \textbf{Magnitude} & \textbf{Mitigation} \\
\midrule
Position bias & Systematic preference for first or second option & 5--15\% & Randomize order \\
\addlinespace
Verbosity bias & Preferring longer responses regardless of quality & 10--20\% & Length-normalize \\
\addlinespace
Self-preference & Preferring outputs from same model family & 10--25\% & Use different judge \\
\addlinespace
Style bias & Preferring confident tone even when wrong & Variable & Rubric-based scoring \\
\addlinespace
Instruction leakage & Rewarding rubric-hacking & Variable & Blind to criteria \\
\bottomrule
\end{tabular}
\end{table}

\paragraph{Position Bias.} LLMs systematically favor the first or second position when comparing two outputs~\citep{wang2024large, zheng2023judging}. In pairwise comparisons, GPT-4 shows approximately 10\% preference for the first position across diverse tasks. Mitigation requires randomizing presentation order and averaging across both orderings.

\paragraph{Verbosity Bias.} Longer responses receive higher ratings regardless of actual quality~\citep{zheng2023judging}. This is particularly problematic when comparing a concise, correct answer against a verbose but less accurate one. Mitigation includes explicit rubric criteria penalizing unnecessary length or length-normalizing scores.

\paragraph{Self-Preference Bias.} LLMs prefer outputs generated by themselves or by similar models~\citep{panickssery2024llm}. When GPT-4 evaluates GPT-4 outputs against Claude outputs, it shows measurable preference for GPT-4 responses. Always use a judge model from a different family than the system being evaluated.

\paragraph{Style Bias.} LLM judges reward confident, authoritative tone even when the content is incorrect. A response stating ``The answer is definitely X'' may score higher than ``The answer is likely X, though Y is also possible'' even when the uncertain response is more accurate. Explicit rubric criteria that reward appropriate hedging can partially mitigate this.

\paragraph{Instruction Leakage (Rubric-Hacking).} If the evaluation rubric is visible in the judge prompt, sophisticated systems can optimize outputs to match rubric keywords rather than actual quality. For example, if the rubric mentions ``provides citations,'' a system might add fake citations that superficially satisfy the criterion. Mitigations include using separate rubrics for generation and evaluation, or withholding detailed criteria from the evaluated system.

\begin{warningbox}
Never use LLM-as-judge as the sole arbiter for high-stakes decisions. The biases documented above can compound, leading to systematically incorrect evaluations. Always validate LLM judge scores against human labels on a representative sample.
\end{warningbox}

\subsection{Good vs.\ Bad Judge Prompts}

The quality of evaluation depends heavily on prompt design.

\begin{badexample}[Vague Judge Prompt]
\begin{lstlisting}[basicstyle=\small\ttfamily]
Is Response A or Response B better?
Just say A or B.
\end{lstlisting}
\textbf{Problems:} No criteria defined; no reasoning required; prone to position bias.
\end{badexample}

\begin{goodexample}[Effective Judge Prompt]
\begin{lstlisting}[basicstyle=\small\ttfamily]
You are an expert evaluator. Compare the two responses below.

[Question]
{question}

[Response A]
{response_a}

[Response B]
{response_b}

Evaluate on these criteria (1-5 each):
1. Accuracy: Are all facts correct?
2. Completeness: Does it fully answer the question?
3. Clarity: Is it well-organized and easy to understand?
4. Conciseness: Is it appropriately brief without 
   unnecessary filler?

First, analyze each response step by step.
Then provide scores in this JSON format:
{
  "analysis_a": "...",
  "analysis_b": "...",
  "scores_a": {"accuracy": X, "completeness": X, ...},
  "scores_b": {"accuracy": X, "completeness": X, ...},
  "winner": "A" or "B" or "tie",
  "reasoning": "..."
}
\end{lstlisting}
\textbf{Strengths:} Explicit criteria; chain-of-thought reasoning; structured output; justification required.
\end{goodexample}

\subsection{Implementation Example}

The following implementation demonstrates position bias mitigation through order randomization:

\begin{lstlisting}[basicstyle=\small\ttfamily,language=Python]
import json

def llm_judge_pairwise(
    question: str,
    response_a: str, 
    response_b: str,
    judge_model: str = "gpt-4"
) -> dict:
    """Evaluate two responses using an LLM judge."""
    
    # Randomize order to mitigate position bias
    if random.random() < 0.5:
        first, second = response_a, response_b
        order = "original"
    else:
        first, second = response_b, response_a
        order = "swapped"
    
    prompt = f"""
    You are an expert evaluator. 
    
    [Question] {question}
    [Response 1] {first}
    [Response 2] {second}
    
    Which response is better? Analyze step by step,
    then output JSON with your verdict.
    """
    
    result = call_llm(judge_model, prompt)
    verdict = json.loads(result)
    
    # Correct for order swapping
    if order == "swapped":
        verdict["winner"] = swap_winner(verdict["winner"])
    
    return verdict
\end{lstlisting}

\subsection{Multi-Judge Aggregation}

Using multiple judge models improves reliability by reducing the impact of model-specific biases:

\begin{lstlisting}[basicstyle=\small\ttfamily,language=Python]
def multi_judge_evaluation(question, response_a, response_b):
    """Use multiple judges and aggregate results."""
    judges = ["gpt-4", "claude-3-opus", "gemini-pro"]
    verdicts = []
    
    for judge in judges:
        verdict = llm_judge_pairwise(
            question, response_a, response_b, judge
        )
        verdicts.append(verdict["winner"])
    
    # Majority vote
    from collections import Counter
    vote_counts = Counter(verdicts)
    winner = vote_counts.most_common(1)[0][0]
    
    # Flag disagreements for human review
    agreement = vote_counts[winner] / len(judges)
    needs_human_review = agreement < 0.67
    
    return {
        "winner": winner,
        "agreement": agreement,
        "needs_human_review": needs_human_review
    }
\end{lstlisting}

\subsection{Comparison: LLM Judges vs.\ Human Evaluation}

Table~\ref{tab:llm-vs-human} compares the trade-offs between LLM judges and human evaluation.

\begin{table}[h]
\centering
\caption{LLM judges vs.\ human evaluation trade-offs.}
\label{tab:llm-vs-human}
\begin{tabular}{@{}lcc@{}}
\toprule
\textbf{Factor} & \textbf{LLM Judge} & \textbf{Human} \\
\midrule
Scale               & High       & Low \\
Cost                & Low        & High \\
Speed               & Fast       & Slow \\
Specialized domains & Variable   & High (experts) \\
Bias awareness      & Limited    & High \\
Subjective nuance   & Moderate   & High \\
Explainability      & Moderate   & High \\
\bottomrule
\end{tabular}
\end{table}

\subsection{Protocol for Robust Judge Evaluation}

To mitigate the biases documented above, we recommend a standardized mitigation protocol rather than ad-hoc prompting.

\begin{itemize}
    \item \textbf{Position Bias Mitigation.} Run every pairwise comparison twice, swapping the order (A vs B, then B vs A). If the judge's preference flips with the order, declare a tie. This double-pass consistency check catches approximately 90\% of sensitivity errors.
    
    \item \textbf{Length Normalization.} To combat verbosity bias, instruct the judge to penalize unnecessary length, or truncate responses to the length of the shorter answer plus 20\% before judging.
    
    \item \textbf{Rubric-Conditioned Scoring.} Instead of asking "Which is better?", ask "Which response better satisfies criteria X, Y, and Z?". Component-level scoring reduces the impact of confident but incorrect answers.
    
    \item \textbf{Reference-Guided Grading.} Providing a gold reference answer grounds the evaluation in truth rather than plausibility, significantly reducing hallucination blind spots.
\end{itemize}

\subsection{Recommended Guardrails}

To deploy LLM judges safely, follow a strict protocol. Never use self-evaluation; a model should not judge its own inputs. Validate judge scores against human labels on a representative sample, targeting a correlation coefficient greater than 0.7. Use judges primarily for screening or regression testing, not as the final certifier for high-stakes decisions. Finally, re-validate calibration whenever the provider updates the model, as judge behavior can drift.

